\begin{table}[H]
  \centering
  \footnotesize
  \begin{tabularx}{\textwidth}{@{}
      >{\raggedright\arraybackslash}p{0.24\textwidth}
      >{\raggedright\arraybackslash}X
      >{\raggedright\arraybackslash}p{0.12\textwidth}@{}}
    \toprule
    Name & One-line claim & Substrate cite \\
    \midrule
    Duality-Fixity Theorem (\Fref{F14}) &
    For a readout that is strongly measurable and square-integrable when the weight is a measure, and a positive cone satisfying (C1)--(C3) (monotone trace, positive trace-zero elements vanish, trace identity), domination by bounds \(B_n\) with \(\inf_n\operatorname{tr}B_n\le0\) forces the anti-invariant
    ledger to vanish and, for a separating readout, confines visible mass
    to \(\Fixinv\) up to a null set; such bounds exist exactly when \(\psi^-=0\) \(\mu\)-almost everywhere.  &
    \PaperSDTC, \PaperRH \\
    Selection-Obstruction Theorem (\Fref{F15a}) &
    For the orbit quotient with its induced action, an equivariant selector
    takes only fixed points; selectedness descends through the quotient exactly
    when its obstruction is empty.  &
    \textemdash \\
    Spontaneous Symmetry Breaking (\Fref{F15b}) &
    A branch-selection record exists exactly when some ground state is
    not invariant; no symmetric selector produces its branch.  &
    \PaperCSL\ (parallel) \\
    Anomaly Theorem (\Fref{F40}) &
    For a claimed symmetry not recorded as explicitly broken, an anomaly
    is the failure of action descent, domain descent or readout
    preservation; for a surjective layer quotient, a descended action is a
    group action (it satisfies the group law and the identity acts
    trivially).  &
    \textemdash \\
    Duality-Equivalence Theorem (\Fref{F41}) &
    In a role-preserving duality equivalence, the inverse also
    intertwines the routes and \(r_B=\rho\circ r_A\circ\bar G\).  &
    \textemdash \\
    \bottomrule
  \end{tabularx}
  \caption{Symmetry at a glance: the five laws of \Cref{sec:symmetry}. Duality Fixity has restricted alignment.}
  \label{tab:axis-symmetry-summary}
\end{table}
